\section{Future Work}
\label{sec:discussion}

\system could allow the user to provide optional constraints on the layout
(either relative or absolute) through pragmas.
A relative constraint would allow the user to specify if a field \lstinline|A| comes immediately after
field \lstinline|B|. An absolute constraint would specify an exact index in the layout for a field.
Such pragmas may be useful if the user requires a specific configuration of a data
type for external reasons or has information about performance bottlenecks.


Although the performance optimization is currently statically driven, there are many avenues for future
improvement. For instance, we can get better edges weights for the access graphs
using dynamic profiling techniques. The profiling can be quite detailed, for
instance, which branch in a function is more likely, which function takes the
most time overall in a global setting (the optimization would bias the layout
towards that function), how does a particular global layout affect the
performance in case of a pipeline of functions.

We could also look at a scenario where we optimize each function locally and 
use ``shim'' functions that copy one layout to another (the one required by the
next function in the pipeline). Although the cost of copying may be high, it warrants
further investigation. Areas of improvement purely on the implementation side include 
optimizing whether \system dereferences a pointer to get to a field or uses the end-witness
information as mentioned in section~\ref{sec:overview}. Lesser pointer dereferencing 
can lower instruction counts and impact performance positively. We would also like to 
optimize the solver times as mentioned in~\ref{subsec:heuristic-tradeoffs}.

We envision that the \textit{structure of arrays} effect that we discovered may help 
with optimizations such as vectorization, where the performance can benefit significantly
if the same datatype is close together in memory. Regardless, through the case studies, we see 
that \system shows promise in optimizing the layout of datatypes and may open up 
the optimization space for other complex optimizations such as vectorization.
